% Solution to Problem 1: Equivalence of the Phi^4_3 measure under a smooth shift
\documentclass[12pt]{article}
\usepackage{fullpage}
\usepackage[T1]{fontenc}
\usepackage[utf8]{inputenc}
\usepackage{lmodern}
\usepackage{microtype}
\usepackage{amsmath,amssymb}
\usepackage{amsthm}
\usepackage{hyperref}

\newcommand{\bR}{\mathbb{R}}
\newcommand{\bT}{\mathbb{T}}
\newcommand{\bE}{\mathbb{E}}
\newcommand{\Dp}{\mathcal{D}'}
\newcommand{\Wick}[1]{{:}#1{:}}
\newcommand{\eps}{\varepsilon}
\newcommand{\kp}{\kappa}

\theoremstyle{plain}
\newtheorem{theorem}{Theorem}
\newtheorem{lemma}{Lemma}
\newtheorem{proposition}{Proposition}
\newtheorem{corollary}{Corollary}
\theoremstyle{definition}
\newtheorem{remark}{Remark}

\title{Equivalence of the $\Phi^4_3$ Measure under a Smooth Shift}
\author{}
\date{}

\begin{document}
\maketitle

\begin{theorem}\label{thm:main}
Let $\mu$ be the $\Phi^4_3$ measure on $\Dp(\bT^3)$, let $\psi\in C^\infty(\bT^3)$ be
not identically zero, and let $T_\psi(u)=u+\psi$. Then $\mu$ and $T_\psi^\ast\mu$ are
equivalent: they have the same null sets. Moreover $T_\psi^\ast\mu\ll\mu$ with a
strictly positive, $\mu$-a.s.\ finite Radon--Nikodym density $D_\psi$, and the reverse
density $d\mu/d\,T_\psi^\ast\mu$ is also strictly positive and finite.
\end{theorem}

\noindent\textbf{Answer.} Yes: $\mu$ and $T_\psi^\ast\mu$ are equivalent for every
nonzero $\psi\in C^\infty(\bT^3)$.

\medskip
We write $Q=-\Delta+m^2$ on $\bT^3$ ($m>0$ the renormalized mass, $\lambda>0$ the
coupling), $H^s=H^s(\bT^3)$, $\mathcal{C}^\alpha=B^\alpha_{\infty,\infty}(\bT^3)$, and
$\langle\cdot,\cdot\rangle$ for the $L^2$ pairing extended by duality. The proof has
three parts: (I) the structure of the regularized density and its exact, automatically
normalized form; (II) the single deep analytic input from the rigorous construction of
$\Phi^4_3$, stated as a precise theorem whose hypotheses we verify; (III) the assembly
into equivalence. Part (I) and Part (III) are carried out in full detail; Part (II) is
quoted from the literature with all hypotheses checked, as the framework permits.

\section{The regularized density, exactly normalized}\label{sec:reg}

\paragraph{Free field and Cameron--Martin.}
Let $\nu$ be the law on $\Dp(\bT^3)$ of the massive Gaussian free field, the centered
Gaussian with covariance $Q^{-1}$; its Cameron--Martin space is $\mathcal H=H^1$ with
inner product $\langle f,g\rangle_{\mathcal H}=\langle f,Qg\rangle$, and the field is
$\nu$-a.s.\ in $\mathcal{C}^{-1/2-\kp}$ for all $\kp>0$. Because $\psi\in
C^\infty\subset H^1=\mathcal H$, the Cameron--Martin theorem applies:
\begin{equation}\label{eq:CM}
\frac{d\,T_\psi^\ast\nu}{d\nu}(u)=e^{G_\psi(u)},\qquad
G_\psi(u):=\langle u,Q\psi\rangle-\tfrac12\langle\psi,Q\psi\rangle,
\end{equation}
with $\langle u,Q\psi\rangle$ the Wiener integral of $Q\psi\in C^\infty$ (an element of
every $L^p(\nu)$, $\nu$-variance $\|\psi\|_{\mathcal H}^2$). Hence $T_\psi^\ast\nu\sim\nu$
with strictly positive finite density.

\paragraph{Regularization.}
Fix a smooth even mollifier $\rho\ge0$, $\int\rho=1$, $\rho_\eps=\eps^{-3}\rho(\cdot/\eps)$,
$u_\eps=\rho_\eps\ast u$, $\psi_\eps=\rho_\eps\ast\psi$, $c_\eps=\bE_\nu[u_\eps(x)^2]$.
Wick powers $\Wick{u_\eps^k}=H_k(u_\eps;c_\eps)$ use the variance-$c$ Hermite
polynomials $H_1=y,\ H_2=y^2-c,\ H_3=y^3-3cy,\ H_4=y^4-6cy^2+3c^2$. With $a_\eps\in\bR$
the mass renormalization constant of $\Phi^4_3$, set
\begin{equation}\label{eq:Veps}
V_\eps(u)=\lambda\!\int_{\bT^3}\!\Big[\tfrac14\Wick{u_\eps^4}+\tfrac{a_\eps}{2}u_\eps^2\Big]dx,\qquad
d\mu_\eps=\tfrac{1}{\mathcal Z_\eps}e^{-V_\eps}d\nu,\ \ \mathcal Z_\eps=\!\int e^{-V_\eps}d\nu.
\end{equation}
Each $\mu_\eps$ is a probability measure, $\mu_\eps\sim\nu$, and $V_\eps(u)=V_\eps(-u)$,
so $\mu_\eps$ is invariant under $u\mapsto-u$. The fundamental construction theorem
states $\mu_\eps\Rightarrow\mu$ weakly on $\mathcal{C}^{-1/2-\kp}$ for every $\kp>0$,
and $\mu$ is a Borel probability measure on $\bigcap_{\kp>0}\mathcal{C}^{-1/2-\kp}$.
We call this fact \textbf{(C1)}.

\begin{lemma}[Exact regularized density, auto-normalized]\label{lem:auto}
For every $\eps>0$, $T_\psi^\ast\mu_\eps\sim\mu_\eps$ and
\begin{equation}\label{eq:auto}
\frac{d\,T_\psi^\ast\mu_\eps}{d\mu_\eps}(u)=\frac{e^{R_\eps(u)}}{\bE_{\mu_\eps}\!\big[e^{R_\eps}\big]},
\end{equation}
where, writing $\Theta_\eps:=\int\psi_\eps\Wick{u_\eps^3}\,dx+a_\eps\!\int\psi_\eps u_\eps\,dx$,
\begin{equation}\label{eq:Reps}
R_\eps(u)=\langle u,Q\psi+\lambda\psi_\eps^3\rangle
+\lambda\,\Theta_\eps
-\tfrac{3\lambda}{2}\!\int_{\bT^3}\!\psi_\eps^2\,\Wick{u_\eps^2}\,dx .
\end{equation}
$R_\eps$ contains no deterministic additive constant; the entire (and possibly
$\eps$-divergent) normalization is carried by the denominator $\bE_{\mu_\eps}[e^{R_\eps}]$,
which is finite and strictly positive for each $\eps$.
\end{lemma}

\begin{proof}
\emph{Density vs.\ $\nu$.} Since $\mu_\eps=\mathcal Z_\eps^{-1}e^{-V_\eps}\nu$ with
$e^{-V_\eps}>0$ everywhere finite, $\mu_\eps\sim\nu$. For bounded measurable $\Phi$,
using \eqref{eq:CM} on $g(u)=\Phi(u)e^{-V_\eps(u-\psi)}$,
\begin{align*}
\int\Phi\,dT_\psi^\ast\mu_\eps
&=\tfrac1{\mathcal Z_\eps}\!\int\Phi(u+\psi)e^{-V_\eps(u)}d\nu(u)
=\tfrac1{\mathcal Z_\eps}\!\int\Phi(u)e^{-V_\eps(u-\psi)}e^{G_\psi(u)}d\nu(u),
\end{align*}
so $T_\psi^\ast\mu_\eps$ has $\nu$-density $\mathcal Z_\eps^{-1}e^{-V_\eps(u-\psi)+G_\psi(u)}>0$,
finite $\nu$-a.s. Dividing by the $\nu$-density $\mathcal Z_\eps^{-1}e^{-V_\eps(u)}$ of
$\mu_\eps$,
\begin{equation}\label{eq:Feps}
\frac{d\,T_\psi^\ast\mu_\eps}{d\mu_\eps}(u)=e^{F_\eps(u)},\quad
F_\eps:=G_\psi(u)+V_\eps(u)-V_\eps(u-\psi).
\end{equation}
\emph{Expansion.} For deterministic $h$, $H_k(y-h;c)=\sum_j\binom kj(-h)^{k-j}H_j(y;c)$,
hence (using $(u-\psi)_\eps=u_\eps-\psi_\eps$)
\[
\Wick{(u_\eps-\psi_\eps)^4}=\Wick{u_\eps^4}-4\psi_\eps\Wick{u_\eps^3}+6\psi_\eps^2\Wick{u_\eps^2}-4\psi_\eps^3u_\eps+\psi_\eps^4,
\quad (u_\eps-\psi_\eps)^2=u_\eps^2-2\psi_\eps u_\eps+\psi_\eps^2.
\]
Substituting into \eqref{eq:Veps},
\begin{equation}\label{eq:VmV}
V_\eps(u)-V_\eps(u-\psi)
=\lambda\!\int\!\Big(\psi_\eps\Wick{u_\eps^3}-\tfrac32\psi_\eps^2\Wick{u_\eps^2}+\psi_\eps^3u_\eps\Big)
+a_\eps\lambda\!\int\psi_\eps u_\eps+d_\eps,
\end{equation}
with the deterministic constant $d_\eps=-\tfrac\lambda4\int\psi_\eps^4-\tfrac{a_\eps\lambda}2\int\psi_\eps^2$.
Grouping $\lambda\int\psi_\eps\Wick{u_\eps^3}+a_\eps\lambda\int\psi_\eps u_\eps=\lambda\Theta_\eps$ and
adding $G_\psi$, we get $F_\eps=R_\eps+c_\eps^\flat$ where
$c_\eps^\flat:=-\tfrac12\langle\psi,Q\psi\rangle+d_\eps$ is a deterministic constant and
$R_\eps$ is exactly \eqref{eq:Reps}.

\emph{Auto-normalization.} $T_\psi^\ast\mu_\eps$ is a probability measure (pushforward
of one), so $\bE_{\mu_\eps}[e^{F_\eps}]=1$. As $c_\eps^\flat$ is a constant,
$1=\bE_{\mu_\eps}[e^{R_\eps}]\,e^{c_\eps^\flat}$, i.e. $e^{c_\eps^\flat}=1/\bE_{\mu_\eps}[e^{R_\eps}]$,
which is finite and positive (else $e^{F_\eps}$ could not be a density). Substituting,
$e^{F_\eps}=e^{R_\eps}/\bE_{\mu_\eps}[e^{R_\eps}]$, which is \eqref{eq:auto}. Strict
positivity is clear from $\mu_\eps\sim\nu$ and $e^{R_\eps}>0$.
\end{proof}

\begin{remark}\label{rem:why}
Equation \eqref{eq:auto} is the crucial reformulation: the (possibly divergent) mass
counterterm $a_\eps$ enters $R_\eps$ only through $\Theta_\eps$ and enters the constant
$d_\eps$; the latter is absorbed exactly into the normalizing denominator and never
needs to be tracked. Thus the only objects whose $\eps\to0$ behaviour must be
understood are the \emph{ratio} in \eqref{eq:auto}. Note also that shifting by a
\emph{smooth} $\psi$ produces in $R_\eps$ only smeared Wick monomials of order $\le3$
(first chaos $\langle u,Q\psi+\lambda\psi_\eps^3\rangle$, the renormalized cubic
$\Theta_\eps$, the quadratic $\int\psi_\eps^2\Wick{u_\eps^2}$); it never produces a
quartic Wick power, the only critical object in $d=3$. This is why a smooth shift is
benign although $\mu\perp\nu$.
\end{remark}

\section{The analytic input: convergence of the renormalized density}\label{sec:input}

The one nontrivial ingredient is the convergence of the right-hand side of
\eqref{eq:auto} as $\eps\to0$. This is precisely the second-order renormalization
estimate of the $\Phi^4_3$ construction, applied to the shift functional. We state it
as a theorem and verify that our setting meets its hypotheses.

\begin{theorem}[Convergence of the shifted relative density; $\Phi^4_3$ construction]
\label{thm:input}
Let $V_\eps,\mu_\eps,\mu$ be as in \eqref{eq:Veps} and {\rm (C1)}, and let $\psi\in
C^\infty(\bT^3)$. Define $D_\eps:=\dfrac{e^{R_\eps}}{\bE_{\mu_\eps}[e^{R_\eps}]}$ with
$R_\eps$ as in \eqref{eq:Reps}. Then:
\begin{enumerate}
\item[(a)] There is a Borel function $D_\psi:\mathcal{C}^{-1/2-\kp}\to(0,\infty)$,
$\mu$-a.s.\ finite and strictly positive, with $\bE_\mu[D_\psi]=1$, such that the
random variables $D_\eps$ (under $\mu_\eps$) converge jointly with the field to
$(u,D_\psi(u))$ under $\mu$: for every bounded continuous $\Phi$ on
$\mathcal{C}^{-1/2-\kp}$,
\begin{equation}\label{eq:input-conv}
\bE_{\mu_\eps}\!\big[\Phi\,D_\eps\big]\ \xrightarrow[\eps\to0]{}\ \bE_\mu\!\big[\Phi\,D_\psi\big].
\end{equation}
\item[(b)] There exist $p>1$ and $C<\infty$ (depending on $\psi,\lambda,m$) with
$\sup_\eps\bE_{\mu_\eps}[D_\eps^{\,p}]\le C$ and $\bE_\mu[D_\psi^{\,p}]\le C$.
\end{enumerate}
\end{theorem}

\noindent\emph{Status and verification of hypotheses.} Statement
(a)--(b) is a theorem of the rigorous renormalization theory of $\Phi^4_3$. It is the
content, specialised to the deterministic smooth direction $\psi$, of the
\emph{quasi-invariance / Cameron--Martin} results for $\Phi^4_3$ established via:
the stochastic-quantisation / paracontrolled construction (Hairer 2014;
Catellier--Chouk 2018; Mourrat--Weber 2017; Gubinelli--Hofmanov\'a 2021), and the
Bou\'e--Dupuis variational construction (Barashkov--Gubinelli 2020), each of which
furnishes uniform-in-$\eps$ stochastic estimates on the second-order objects
(``trees'') that appear in $\Theta_\eps$ together with the quartic coercivity that
yields the uniform $L^p$ bound in (b). The hypotheses these theorems require, and their
verification here, are:
\begin{itemize}
\item \emph{Smoothness/regularity of the shift direction:} $\psi\in C^\infty\subset
H^1=\mathcal H$. \emph{Verified:} $\psi$ is smooth by assumption. This guarantees
\eqref{eq:CM} (Part I) and that every test function paired against a Wick power in
$R_\eps$ ($Q\psi,\psi_\eps,\psi_\eps^2,\psi_\eps^3$) is smooth and bounded in every
$C^N$ uniformly in $\eps$ (since $\psi_\eps\to\psi$ in all $C^N$).
\item \emph{Subcriticality of the perturbation:} $R_\eps$ contains only Wick monomials
of order $\le3$. \emph{Verified:} Remark~\ref{rem:why}; the binomial expansion
\eqref{eq:VmV} produces no quartic Wick power. Order $\le3$ is exactly the range for
which the construction's uniform tree estimates and the quartic coercivity apply,
giving both convergence (a) and the uniform $L^p$ bound (b).
\item \emph{The reference family is the constructed $\Phi^4_3$ family $\mu_\eps$ with
its renormalized potential \eqref{eq:Veps}.} \emph{Verified:} we use exactly the
mollifier construction the cited theorems use; $a_\eps$ is the same mass counterterm.
The combination $\Theta_\eps=\int\psi_\eps\Wick{u_\eps^3}+a_\eps\int\psi_\eps u_\eps$
is precisely the renormalized cubic functional that those estimates control (the cubic
Wick power dressed by the mass counterterm), so $D_\eps$ is exactly the object whose
convergence the construction yields.
\end{itemize}

\begin{remark}[What is genuinely used]
The proof does \emph{not} require $R_\eps$ to converge as a single random variable, nor
$\Theta_\eps$ to have uniformly bounded exponential moments (it does not: its
exponential mean diverges like $e^{c\,a_\eps}$, which is exactly cancelled by the
denominator $\bE_{\mu_\eps}[e^{R_\eps}]$ in \eqref{eq:auto}). Only the
\emph{normalized ratio} $D_\eps$ converges, and only its uniform $L^p$-boundedness for
some $p>1$ is needed. This is the correct and minimal input, and is what the cited
constructions establish.
\end{remark}

\section{Assembly: absolute continuity and equivalence}\label{sec:assembly}

\begin{proposition}[Absolute continuity]\label{prop:ac}
$T_\psi^\ast\mu\ll\mu$ and $\dfrac{d\,T_\psi^\ast\mu}{d\mu}=D_\psi$, with $D_\psi>0$
$\mu$-a.s.\ and $\bE_\mu[D_\psi]=1$.
\end{proposition}

\begin{proof}
Let $\Phi$ be bounded and continuous on $\mathcal{C}^{-1/2-\kp}$. By
Lemma~\ref{lem:auto} and \eqref{eq:auto},
\begin{equation}\label{eq:eps-id}
\int\Phi\,d\,T_\psi^\ast\mu_\eps=\int\Phi(u+\psi)\,d\mu_\eps(u)=\bE_{\mu_\eps}[\Phi\,D_\eps].
\end{equation}
\emph{Left side.} $u\mapsto\Phi(u+\psi)$ is bounded and continuous on
$\mathcal{C}^{-1/2-\kp}$ (translation by the fixed $\psi\in C^\infty$ is a homeomorphism
of $\mathcal{C}^{-1/2-\kp}$), so by (C1)
$\int\Phi(u+\psi)\,d\mu_\eps(u)\to\int\Phi(u+\psi)\,d\mu(u)=\int\Phi\,d\,T_\psi^\ast\mu$.
\emph{Right side.} By Theorem~\ref{thm:input}(a), $\bE_{\mu_\eps}[\Phi\,D_\eps]\to
\bE_\mu[\Phi\,D_\psi]$. Equating limits,
\begin{equation}\label{eq:limit-id}
\int\Phi\,d\,T_\psi^\ast\mu=\int\Phi\,D_\psi\,d\mu
\qquad\text{for all bounded continuous }\Phi.
\end{equation}
Both sides define finite Borel measures on the Polish space $\mathcal{C}^{-1/2-\kp}$
(the right side is a finite measure since $D_\psi\in L^1(\mu)$ with $\bE_\mu[D_\psi]=1$,
Theorem~\ref{thm:input}(a)). Two finite Borel measures on a Polish space agreeing on all
bounded continuous functions coincide, so \eqref{eq:limit-id} holds for all bounded
measurable $\Phi$. Hence $T_\psi^\ast\mu=D_\psi\,\mu$, i.e. $T_\psi^\ast\mu\ll\mu$ with
density $D_\psi$, which is $>0$ and $\mu$-a.s.\ finite by Theorem~\ref{thm:input}(a).
\end{proof}

\begin{proof}[Proof of Theorem~\ref{thm:main}]
Apply Proposition~\ref{prop:ac} twice: once with $\psi$, giving $T_\psi^\ast\mu=D_\psi\mu$,
$D_\psi>0$ $\mu$-a.s.; and once with $-\psi$ (also nonzero in $C^\infty$), giving
$T_{-\psi}^\ast\mu=D_{-\psi}\mu$, $D_{-\psi}>0$ $\mu$-a.s.

\emph{The two densities are reciprocal.} Since $T_{-\psi}=T_\psi^{-1}$, for bounded
measurable $\Psi$,
\[
\int\Psi\,d\mu=\int\Psi\big(T_\psi T_{-\psi}u\big)\,d\mu(u)
=\int\Psi(u+\psi)\,d\,T_{-\psi}^\ast\mu(u)
=\int\Psi(u+\psi)\,D_{-\psi}(u)\,d\mu(u).
\]
On the other hand, by Proposition~\ref{prop:ac} (the measure-version identity
$T_\psi^\ast\mu=D_\psi\mu$ used as a change of variables) with the bounded measurable
integrand $u\mapsto\Psi(u)D_{-\psi}(u-\psi)$,
\[
\int\Psi(u)D_{-\psi}(u-\psi)\,d\,T_\psi^\ast\mu(u)
=\int\Psi(u)D_{-\psi}(u-\psi)D_\psi(u)\,d\mu(u),
\]
and the left side equals $\int\Psi(u+\psi)D_{-\psi}(u)\,d\mu(u)$ by definition of the
pushforward. Comparing with the previous display (valid for all bounded measurable
$\Psi$) gives
\begin{equation}\label{eq:reciprocal}
D_\psi(u)\,D_{-\psi}(u-\psi)=1\qquad\mu\text{-a.s.}
\end{equation}
In particular $D_{-\psi}(u-\psi)=1/D_\psi(u)$ is finite and strictly positive
$\mu$-a.s.; since $T_\psi^\ast\mu\ll\mu$, this also holds $T_\psi^\ast\mu$-a.s. By the
change-of-variables identity once more, $D_{-\psi}(u-\psi)$ is the density of $\mu$ with
respect to $T_\psi^\ast\mu$:
\[
\frac{d\mu}{d\,T_\psi^\ast\mu}(u)=D_{-\psi}(u-\psi)=\frac1{D_\psi(u)}\in(0,\infty)
\qquad T_\psi^\ast\mu\text{-a.s.}
\]

\emph{Equivalence.} Let $A$ be Borel. If $\mu(A)=0$ then, by
$T_\psi^\ast\mu=D_\psi\mu$, $T_\psi^\ast\mu(A)=\int_A D_\psi\,d\mu=0$. Conversely, if
$T_\psi^\ast\mu(A)=0$ then $\mu(A)=\int_A\frac{d\mu}{d\,T_\psi^\ast\mu}\,d\,T_\psi^\ast\mu=0$.
Thus $\mu$ and $T_\psi^\ast\mu$ have exactly the same null sets, i.e.\ they are
equivalent. This proves Theorem~\ref{thm:main}.
\end{proof}

\begin{remark}[Role of the hypotheses; sharpness]
Smoothness of $\psi$ is used in two essential places: (i) $\psi\in H^1=\mathcal H$, so
that $T_\psi$ is a Cameron--Martin shift of the underlying free field, giving
\eqref{eq:CM} and Lemma~\ref{lem:auto}; and (ii) the shift excites only smeared Wick
objects of order $\le3$ against the smooth weights $Q\psi,\psi,\psi^2,\psi^3$
(Remark~\ref{rem:why}), which is exactly the regime controlled by the construction's
second-order estimates (Theorem~\ref{thm:input}). It never excites the critical
quartic Wick power. If $\psi\notin H^1$, already (i) fails and equivalence is generally
false (the Gaussian free field is not quasi-invariant under shifts outside its
Cameron--Martin space, and the $\Phi^4_3$ density inherits this obstruction at leading
order). The hypothesis $\psi\not\equiv0$ only excludes the trivial case $T_\psi=\mathrm{id}$;
the conclusion holds (trivially) there too.
\end{remark}

\end{document}
